\documentclass{exam}
\usepackage{../commonheader}

\discnumber{3}
\title{Python Control, Expressions \& Higher Order Functions}
\date{September 8 -- September 12, 2025}

\begin{document}
\maketitle

\begin{meta}
    \textbf{Example Timeline}
    \begin{itemize}
        \item Tips for Mentors
            \subitem{Note: \textbf{Each problem's "Suggested Time" is not assigned with the 1 hour time for sections in mind.} It is more of a note of how long each individual problem may take. That is, if you added all the "Suggested Times" in the coding section, it could very well be greater than 1 hour (especially in more complicated worksheets in the future).}
            \subitem{\textbf{However, the time in brackets will add up to ~50 minutes, and is intended to be a suggested outline for your section.}}

        \item Intros \& Logstics [6 min]
            \subitem{Introduce yourself! Go over how CSM works/answer any logistical questions.}
        \item Mini-lecture abt control/functions [5 min]  
        \item Q1: WWPD: Potpourri [6 min]
            \subitem{Feel free to skip parts!}
        \item Q2: Handle Overflow [7 min]
        \item Q3: FizzBuzz [7 min]
        \item Mini-lecture abt higher order functions [10 min] 
        \item Q4: checksum [10 min]
        \item If extra time: Q5: curry.forever.alt
        \item How to be successful in 61A [5 min]
            \subitem{Give any tips you want!}
    \end{itemize}
\end{meta}

\section{Intro to Python}
\begin{questions}
    \subimport{../../topics/basics/}{potpurri.tex}
    \begin{questionmeta}
        Suggested Time: 6 min
        \begin{itemize}
            \item Again, you can probably skip a lot of these as soon as your students get the idea.
            \item Your students might be new to solving "compound" questions like the last one; Consider giving them tips on how to break down such problems.
        \end{itemize}
    \end{questionmeta}
\end{questions}


\section{Control}
\begin{questions}
    \subimport{../../topics/control/easy/}{handle-overflow.tex}
    \begin{questionmeta}
        Suggested Time: 7 min; Difficulty: Easy
        \begin{itemize}
            \item Make sure students understand the difference between using else and elif in these types of problems.
        \end{itemize}
    \end{questionmeta}

    \subimport{../../topics/control/medium/}{fizzbuzz.tex}
    \begin{questionmeta}
        Suggested Time: 7 min; Difficulty: Medium
        \begin{itemize}
            \item This question has \lstinline{while} loops! You may want to discuss the difference between \lstinline{for} and \lstinline{while} loops. 
            \item Also make sure your students understand that print returns None [as seen in the final doctest] and that it is valid for a Python function to not have a return statement.
        \end{itemize}
    \end{questionmeta}
\end{questions}

\section{Higher Order Functions}
\begin{questions}
    \subimport{../../topics/hof/medium/}{check_sum.tex}
    \begin{questionmeta}
        Suggested Time: 10 min; Difficulty: Medium
        \begin{itemize}
            \item This question has a lot of blank lines and also combines multiple topics, like while loops, that students have been learning
            \item Focus on why we return check, and not a boolean! 
        \end{itemize}
    \end{questionmeta}

    \subimport{../../topics/hof/medium/}{curry_forever_alt.tex}
    \begin{questionmeta}
        Suggested Time: 10 min; Difficulty: Medium
        \begin{itemize}
            \item Make sure to show how we can use argnum to terminate the function at the end
            \item Also focus on how f is being passed as a parameter to curry forever, and why we need to do f(x, base)
        \end{itemize}
    \end{questionmeta}

\end{questions}

\end{document}